\documentclass{article}
\usepackage{amsthm}
\usepackage{comment}
\usepackage{amsmath}
\usepackage{amsfonts}
\usepackage{sectsty}
\usepackage{hyperref}
\usepackage{fancyhdr}
\usepackage[top=1in,bottom=1in,left=1in,right=1in]{geometry}
\usepackage{float}
\usepackage{graphicx}
\usepackage{tabularx}
\usepackage{mathtools}

\newtheorem{theorem}{Theorem}
\newtheorem{corollary}{Corollary}[theorem]

\title{Lecture 5 - Properties of Fields}
\author{Matthew Farah, \href{mailto:farahm11@mcmaster.ca}{$\langle$farahm11@mcmaster.ca$\rangle$}, 400468251}
\date{\today}

\hypersetup{
        colorlinks=true,
        linkcolor=blue,
        filecolor=magenta,
        urlcolor=blue,
        citecolor=blue,
	anchorcolor=blue
}

\fancyhead{}
\fancyhead[L]{Matthew Farah, \href{mailto:farahm11@mcmaster.ca}{$\langle$farahm11@mcmaster.ca$\rangle$}, 400468251}
\fancyhead[R]{Lecture 5 - Properties of Fields}

\newenvironment{directsubsubs}{
  \makeatletter
  \renewcommand\thesubsubsection{\thesection.\Roman{subsubsection}}
  \makeatother
}{
  \makeatletter
  \renewcommand\thesubsubsection{\thesubsection.\Roman{subsubsection}}
  \makeatother
}

\renewcommand{\thesubsection}{\thesection.\alph{subsection}}
\renewcommand{\thesubsubsection}{\thesection.\alph{subsection}.\Roman{subsubsection}}
\makeatletter

\sectionfont{\fontsize{12}{12}\bfseries}
\subsectionfont{\fontsize{10}{12}\normalfont}
\subsubsectionfont{\fontsize{10}{12}\normalfont}

\begin{document}

\maketitle
\pagestyle{fancy}

From our prior definition of a field, we can construct a number of useful theorems.

\begin{theorem}
	The addition-cancellation property: $x + y = x + z \implies y = z$.
\end{theorem}

\begin{proof}
	\begin{align*}
		y &= 0 + y\\
		&= (-x + x) + y \\
		&= -x + (x + y) \\
		&= -x + (x + z) \\
		&= (-x + x) + z \\
		&= 0 + z \\
		&= z
	\end{align*}
\end{proof}

\begin{corollary}
	The additive inverse of any element in a field is unique
\end{corollary}

\begin{proof}
	For the purposes of contradiction, assume that for $x \in \mathbb{F}$, there exists two distinct additive inverses $-x_1, -x_2 \in \mathbb{F}, \; -x_1 -x_2 = 0$ of $x$. Then, we have:
	
	\begin{align*}
		-x_1 + x &= 0 \\
		&= -x_2 + x
	\end{align*}

	Which, by the addition-cancellation property, implies that $-x_1 = -x_2$, and thus, by contradiction, the additive inverse of an element in $\mathbb{F}$ is always unique
\end{proof}

\begin{proof}
	We can prove this by considering the product of any number $x \in \mathbb{F}$ with 0.

	\begin{align*}
		0 \cdot x + 0 \cdot y &= (0 + 0) \cdot x \\
		&= 0 + \cdot x
	\end{align*}

	Therefore, by the addition-cancellation property, $0 \cdot x = 0 \; \forall x \in \mathbb{F}$.
\end{proof}

\section[Ex. ~\thesection]{Example: Prove that the product of anything with zero is always zero}


\end{document}
